\documentclass[11pt]{article}

\usepackage[margin=1in]{geometry}
\usepackage{amsmath}
\usepackage{amssymb}
\usepackage{booktabs}
\usepackage{enumitem}
\usepackage{hyperref}

\title{Posterior-Window Batched Monotone Tracing}
\author{}
\date{}

\begin{document}

\maketitle

\section{Goal}

The method estimates the fidelity-$0.5$ boundary in randomized benchmarking
configuration space
\[
    (n,r) =
    (\texttt{n\_gates}, \texttt{ratio\_2\_qb\_gates}),
\]
where $n$ is the total gate count and $r$ is the fraction of two-qubit gates.
For each fixed ratio $r$, the target boundary value is
\[
    n_*(r) \quad \text{such that} \quad p(n_*(r),r)=0.5,
\]
where $p(n,r)$ is the survival fidelity/probability estimated by randomized
benchmarking circuits.

The algorithm traces the boundary from high two-qubit ratio to low two-qubit
ratio. At each ratio it probes a small batch of gate counts. The batch window
is chosen from a particle posterior over a simple inverse decay model. All
measured configurations update this posterior; accepted local crossings are
recorded for plotting and diagnostics, but the posterior update does not
depend on accepting a crossing.

\section{Configuration and Data}

The experiment uses the same RMB data model as the other crossing experiments.
Each probed configuration is an \texttt{RMBConfig} with fixed number of qubits,
native gate set, a one-qubit gate count, and a two-qubit gate count. The helper
\texttt{settings.make\_config(n,r)} maps total gates and two-qubit ratio to a
valid circuit configuration by rounding to realizable even gate counts.

For each measured configuration $x_i=(n_i,r_i)$, the data store contains a
Bayesian estimator with Bernoulli counts
\[
    y_i = \#\{\text{successful circuits at } x_i\},
    \qquad
    z_i = \#\{\text{failed circuits at } x_i\}.
\]
The likelihood used by this method works directly with these counts.

The main settings controlling the method are summarized below.

\begin{center}
\begin{tabular}{p{0.48\textwidth}p{0.42\textwidth}}
\toprule
Setting & Meaning \\
\midrule
\texttt{ratio\_step} & Fixed ratio step if adaptive stepping is disabled. \\
\texttt{adaptive\_ratio\_step} & Whether to adapt the ratio step to the remaining budget. \\
\texttt{min\_ratio\_step}, \texttt{max\_ratio\_step} & Bounds on adaptive ratio steps. \\
\texttt{target\_hqc\_per\_ratio} & Fallback budget scale used before posterior scoring is available. \\
\texttt{adaptive\_ratio\_candidates} & Number of candidate next ratios scored by acquisition value. \\
\texttt{adaptive\_small\_step\_penalty\_power} & Penalty exponent discouraging unnecessarily small adaptive steps. \\
\texttt{adaptive\_cost\_power} & Exponent controlling how strongly expected HQC cost penalizes a candidate. \\
\texttt{adaptive\_coverage\_bonus\_weight} & Weight of a small progress/coverage bonus in adaptive ratio selection. \\
\texttt{initial\_one\_q\_pauli\_error} & Prior one-qubit Pauli error estimate. \\
\texttt{initial\_two\_q\_pauli\_error} & Prior two-qubit Pauli error estimate. \\
\texttt{initial\_error\_relative\_uncertainty} & Default relative uncertainty on both error estimates. \\
\texttt{initial\_window\_confidence} & Confidence interval used for the first analytic window. \\
\texttt{initial\_window\_min\_width} & Minimum gate width of the first window. \\
\texttt{initial\_batch\_points}, \texttt{initial\_batch\_shots} & First-batch size and shots per point. \\
\texttt{trace\_batch\_points}, \texttt{trace\_batch\_shots} & Later-batch size and shots per point. \\
\texttt{max\_hqc\_per\_ratio} & Per-ratio budget cap after data exist. \\
\texttt{particle\_count} & Number of prior particles for the posterior approximation. \\
\texttt{particle\_sharpness} & Prior median of the logistic transition sharpness. \\
\texttt{particle\_sharpness\_log\_std} & Log prior width of the sharpness. \\
\texttt{particle\_curvature\_std\_fraction} & Prior width of the signed curvature as a fraction of the slope prior. \\
\texttt{posterior\_window\_quantiles} & Posterior gate-count quantiles defining each trace window. \\
\texttt{posterior\_probe\_quantiles} & Posterior gate-count quantiles used as probe locations. \\
\texttt{posterior\_window\_padding\_fraction} & Padding added to the posterior window. \\
\texttt{posterior\_window\_min\_width} & Minimum gate width of posterior windows. \\
\texttt{posterior\_window\_max\_width\_fraction} & Optional cap on posterior window width. \\
\texttt{posterior\_min\_configs} & Minimum measured configurations before posterior windows are used. \\
\bottomrule
\end{tabular}
\end{center}

\section{Ratio Schedule}

The trace is always top-down. Let the ratio bounds be
\[
    r_{\min}, r_{\max} = \texttt{settings.ratio\_bounds}.
\]
The first ratio is $r_{\max}$ and the trace proceeds toward $r_{\min}$.

If \texttt{adaptive\_ratio\_step} is false, the schedule is the inclusive
fixed-step sequence
\[
    r_0=r_{\max},\qquad
    r_{k+1}=\max(r_{\min}, r_k-\texttt{ratio\_step}).
\]

If \texttt{adaptive\_ratio\_step} is true but fewer than
\texttt{posterior\_min\_configs} measured configurations exist, the method
uses a fallback budget-spreading step. With remaining HQC budget $B$,
\[
    A = \max\left(1,\left\lfloor
        \frac{B}{\texttt{target\_hqc\_per\_ratio}}
    \right\rfloor\right),
\]
\[
    \Delta r =
    \operatorname{clip}
    \left(
        \frac{|r_{\min}-r_k|}{A},
        \texttt{min\_ratio\_step},
        \texttt{max\_ratio\_step}
    \right),
\]
\[
    r_{k+1}=\max(r_{\min}, r_k-\Delta r).
\]
This fallback is only used until a posterior fit can be formed.

Once at least \texttt{posterior\_min\_configs} measured configurations exist,
the next ratio is chosen by scoring candidate downward steps. Let
\[
    S_{\min}=\min(\texttt{min\_ratio\_step}, r_k-r_{\min}),
    \qquad
    S_{\max}=\min(\texttt{max\_ratio\_step}, r_k-r_{\min}).
\]
The method forms \texttt{adaptive\_ratio\_candidates} evenly spaced candidate
steps
\[
    S_c \in [S_{\min},S_{\max}],
    \qquad
    r_c = r_k-S_c.
\]
For each candidate ratio $r_c$, the method first builds the posterior window
and probe gate counts that would be submitted there. Let those candidate probe
gate counts be $h_1,\ldots,h_K$. For each posterior particle $\theta_j$, define
the model probability at probe $k$ by
\[
    p_{jk} = \pi_{\theta_j}(h_k,r_c).
\]
The posterior predictive mean and particle variance at the probe are
\[
    \bar p_k = \sum_j w_j p_{jk},
    \qquad
    V_k = \sum_j w_j(p_{jk}-\bar p_k)^2.
\]
The acquisition value of the candidate batch is
\[
    I(r_c)
    =
    \sum_{k=1}^K
    V_k\,4\bar p_k(1-\bar p_k).
\]
The variance term rewards probes where plausible boundary curves disagree.
The factor $4\bar p_k(1-\bar p_k)$ emphasizes probes whose predicted outcome is
near the threshold, where a Bernoulli observation is most informative. The
method also estimates the stitched HQC cost $C(r_c)$ of submitting the candidate
batch using the same batch-cost accounting as the real submission.

The candidate score is
\[
    \operatorname{score}(r_c)
    =
    \frac{I(r_c)}{C(r_c)^\eta}
    \left(
        \frac{S_c}{S_{\max}}
    \right)^\gamma
    +
    \lambda
    \frac{S_c}{S_{\max}},
\]
where
\[
    \begin{aligned}
    \gamma &= \texttt{adaptive\_small\_step\_penalty\_power},\\
    \eta &= \texttt{adaptive\_cost\_power},\\
    \lambda &= \texttt{adaptive\_coverage\_bonus\_weight}.
    \end{aligned}
\]
The next ratio is the candidate with the largest score. This selects batches
expected to discriminate between plausible fitted curves, uses cost as a
soft penalty rather than a hard divisor, and adds a small bonus for making
bounded progress through the top-down trace.

\section{Boundary Model}

The fitted boundary model has inverse form
\[
    n_\theta(r) = \frac{1}{q+s r+c r(1-r)},
\]
where
\[
    \theta=(q,s,c,\kappa).
\]
Here $q>0$ is the one-qubit/intercept decay parameter, $s>0$ is the
two-qubit-ratio slope parameter, $c$ is a signed curvature parameter, and
$\kappa>0$ is a transition sharpness. The curvature basis $r(1-r)$ allows the
fit to bend in the middle of the ratio interval while leaving the endpoint
interpretation of $q$ and $s$ relatively stable.

The probability that a configuration $(n,r)$ is above the fidelity-$0.5$
threshold is modeled by a logistic transition around $n_\theta(r)$:
\[
    \pi_\theta(n,r)
    =
    \sigma\left(
        \kappa \frac{n_\theta(r)-n}{n_\theta(r)}
    \right),
\]
where
\[
    \sigma(t)=\frac{1}{1+\exp(-t)}.
\]
Thus, when $n<n_\theta(r)$, the model predicts probability above $0.5$ greater
than one half; when $n>n_\theta(r)$, it predicts probability below one half.

\section{Initial Analytic Prior}

The particle prior is centered on an analytic line derived from the configured
initial Pauli error guesses. Let
\[
    e_1 = \texttt{initial\_one\_q\_pauli\_error},
    \qquad
    e_2 = \texttt{initial\_two\_q\_pauli\_error}.
\]
Let the default base errors be
\[
    e_{1,0}=2.5\times 10^{-5},
    \qquad
    e_{2,0}=7.9\times 10^{-4}.
\]
The prior center for the inverse boundary is
\[
    q_0 =
    \frac{\texttt{REFERENCE\_OFFSET}}{\log 2}
    \frac{e_1}{e_{1,0}},
    \qquad
    s_0 =
    \frac{\texttt{REFERENCE\_SLOPE}}{\log 2}
    \frac{e_2}{e_{2,0}}.
\]
This gives the prior boundary
\[
    n_0(r)=\frac{1}{q_0+s_0 r}.
\]
The curvature prior is centered at zero, so the analytic prior remains the
linear inverse boundary before data are observed.

\subsection{Initial Window}

The first batch is not chosen from the particle posterior, because no
measurements are available. Instead, it uses an analytic uncertainty interval
around the initial error guesses.

Let the relative uncertainties be
\[
    u_1 =
    \begin{cases}
    \texttt{initial\_one\_q\_error\_relative\_uncertainty},
    & \text{if not None},\\
    \texttt{initial\_error\_relative\_uncertainty},
    & \text{otherwise},
    \end{cases}
\]
\[
    u_2 =
    \begin{cases}
    \texttt{initial\_two\_q\_error\_relative\_uncertainty},
    & \text{if not None},\\
    \texttt{initial\_error\_relative\_uncertainty},
    & \text{otherwise}.
    \end{cases}
\]
For confidence level $c=\texttt{initial\_window\_confidence}$, define
\[
    z_c = \Phi^{-1}\left(\frac{1+c}{2}\right),
\]
where $\Phi$ is the standard normal CDF. The lower and upper multiplicative
factors are
\[
    a_1^-=\max(10^{-9},1-z_c u_1), \quad
    a_1^+=1+z_c u_1,
\]
\[
    a_2^-=\max(10^{-9},1-z_c u_2), \quad
    a_2^+=1+z_c u_2.
\]
The method evaluates the analytic crossing at three error pairs:
\[
    (e_1,e_2), \qquad (a_1^+e_1,a_2^+e_2),
    \qquad (a_1^-e_1,a_2^-e_2).
\]
The high-error pair gives the low-gate side of the window and the low-error
pair gives the high-gate side. If the resulting width is less than
\texttt{initial\_window\_min\_width}, the window is widened symmetrically
around the central analytic prediction.

Finally the window is clamped to \texttt{n\_gates\_bounds} and rounded to even
gate counts. If the analytic window cannot be formed, the fallback first
window is the full gate-count range.

\section{Particle Prior}

At every posterior update, the method creates a fresh set of
$M=\texttt{particle\_count}$ prior particles. For particle $j$, it samples
\[
    \log q_j \sim
    \mathcal{N}\left(
        \log q_0,\;
        \sigma_q^2
    \right),
\]
\[
    \log s_j \sim
    \mathcal{N}\left(
        \log s_0,\;
        \sigma_s^2
    \right),
\]
\[
    \log \kappa_j \sim
    \mathcal{N}\left(
        \log \kappa_0,\;
        \sigma_\kappa^2
    \right),
\]
and the signed curvature
\[
    c_j \sim \mathcal{N}(0,\sigma_c^2),
\]
where
\[
    \kappa_0=\texttt{particle\_sharpness},
    \qquad
    \sigma_\kappa =
    \max(0.05,\texttt{particle\_sharpness\_log\_std}),
\]
and
\[
    \sigma_q=\max(0.05,\log(1+u_1)),
    \qquad
    \sigma_s=\max(0.05,\log(1+u_2)).
\]
The curvature width is
\[
    \sigma_c =
    \texttt{particle\_curvature\_std\_fraction}\cdot s_0.
\]
The implementation clips sampled values to numerical ranges
\[
    10^{-8}\le q_j\le 1,\qquad
    10^{-8}\le s_j\le 1,\qquad
    -1\le c_j\le 1,\qquad
    0.05\le \kappa_j\le 100.
\]

The random seed for a posterior update is derived deterministically from the
run seed, the ratio, and the update index:
\[
    \texttt{seed}_{r,t}
    =
    \texttt{rng\_seed}
    + \texttt{particle\_seed\_offset}
    + \operatorname{round}(10^6 r)
    + 997t,
\]
unless \texttt{rng\_seed} is \texttt{None}, in which case fresh entropy is
used.

\section{Likelihood and Posterior Update}

Suppose the measured data after some batches are
\[
    \mathcal{D}
    =
    \{(n_i,r_i,y_i,z_i)\}_{i=1}^N,
\]
where $y_i$ is the number of successful circuits and $z_i$ is the number of
failed circuits at configuration $(n_i,r_i)$. For a particle $\theta_j$, define
\[
    p_{ij} = \pi_{\theta_j}(n_i,r_i).
\]
The Bernoulli likelihood of all measured outcomes is
\[
    L_j
    =
    \prod_{i=1}^N
    p_{ij}^{y_i}
    (1-p_{ij})^{z_i}.
\]
Equivalently, the log likelihood is
\[
    \ell_j =
    \sum_{i=1}^N
    \left[
        y_i\log(p_{ij}+\epsilon)
        +
        z_i\log(1-p_{ij}+\epsilon)
    \right],
\]
with $\epsilon=10^{-12}$ for numerical stability.

The particle weights are normalized exponentiated log likelihoods:
\[
    w_j =
    \frac{\exp(\ell_j-\max_k \ell_k)}
         {\sum_{m=1}^M \exp(\ell_m-\max_k \ell_k)}.
\]
If the normalization fails numerically, the method falls back to uniform
weights. The effective sample size printed in progress messages is
\[
    \operatorname{ESS}
    =
    \frac{1}{\sum_{j=1}^M w_j^2}.
\]

This is the complete Bayesian update used by the method. There is no accepted
crossing constraint, bisection state, or special one-sided correction in the
posterior update. A batch where all measured fidelities are too low or too
high still updates the posterior through the Bernoulli likelihood.

\section{Posterior Fit and Posterior Draws}

The plotting and scoring interface expects a fitted boundary and a collection
of boundary samples. The method provides these from the particle posterior.

\subsection{Posterior Median Fit}

If fewer than \texttt{posterior\_min\_configs} measured configurations exist,
no fit is returned. Otherwise the method computes marginal weighted medians:
\[
    \hat q = Q_{0.5}(\{q_j\},\{w_j\}), \qquad
    \hat s = Q_{0.5}(\{s_j\},\{w_j\}),
\]
\[
    \hat c = Q_{0.5}(\{c_j\},\{w_j\}), \qquad
    \hat \kappa = Q_{0.5}(\{\kappa_j\},\{w_j\}),
\]
where $Q_\alpha$ denotes the weighted quantile. The displayed fitted boundary
is
\[
    \hat n(r)=\frac{1}{\hat q+\hat s r+\hat c r(1-r)}.
\]

\subsection{Posterior Draws}

For uncertainty bands, the method samples particle indices with replacement:
\[
    J_m \sim \operatorname{Categorical}(w_1,\ldots,w_M),
    \qquad
    \theta^{(m)}=\theta_{J_m}.
\]
Each sampled parameter vector gives a boundary draw
\[
    n^{(m)}(r)=
    \frac{1}{q^{(m)}+s^{(m)}r+c^{(m)}r(1-r)}.
\]
Although this is exposed through a method named \texttt{boundary\_bootstrap}
for compatibility with the shared plotting code, these are posterior particle
draws rather than bootstrap refits.

\section{Choosing a Posterior Batch Window}

After enough measured configurations exist, the next batch at ratio $r$ is
chosen from the posterior predictive distribution of $n_*(r)$. For each
particle,
\[
    g_j(r)=\frac{1}{q_j+s_j r+c_jr(1-r)}.
\]
Particles with non-finite or non-positive $g_j(r)$ are discarded and the
remaining weights are renormalized. Curvature particles whose denominator is
non-positive at the evaluated ratio are therefore excluded from that window.

Let
\[
    (\alpha_{\mathrm{lo}},\alpha_{\mathrm{hi}})
    =
    \texttt{posterior\_window\_quantiles}.
\]
The raw window is
\[
    L = Q_{\alpha_{\mathrm{lo}}}(\{g_j(r)\},\{w_j\}),
    \qquad
    U = Q_{\alpha_{\mathrm{hi}}}(\{g_j(r)\},\{w_j\}),
\]
with posterior median
\[
    M = Q_{0.5}(\{g_j(r)\},\{w_j\}).
\]
The window is padded by
\[
    P = \max\left(2,\;
        \texttt{posterior\_window\_padding\_fraction}\cdot(U-L)
    \right),
\]
so that
\[
    L \leftarrow L-P, \qquad U \leftarrow U+P.
\]
If the padded width is below \texttt{posterior\_window\_min\_width}, it is
expanded symmetrically around $M$:
\[
    L \leftarrow M-\frac{W_{\min}}{2}, \qquad
    U \leftarrow M+\frac{W_{\min}}{2},
\]
where
\[
    W_{\min}=\max(2,\texttt{posterior\_window\_min\_width}).
\]
If \texttt{posterior\_window\_max\_width\_fraction} is not \texttt{None}, the
width is capped at
\[
    W_{\max} =
    \max\left(
        W_{\min},
        \texttt{posterior\_window\_max\_width\_fraction}\cdot\max(M,1)
    \right).
\]
If $U-L>W_{\max}$, the window is replaced by
\[
    [M-W_{\max}/2,\; M+W_{\max}/2].
\]
The final window is clamped to \texttt{n\_gates\_bounds} and rounded to even
gate counts.

\section{Choosing Probe Points}

Within a posterior window $[L,U]$, probe locations are chosen using
\texttt{posterior\_probe\_quantiles}. If the number of requested quantiles does
not match the requested number of batch points, the method instead uses
equally spaced quantiles between the posterior window quantiles.

For requested probe quantiles $\beta_1,\ldots,\beta_K$, the raw probe gate
counts are
\[
    h_k =
    2\operatorname{round}
    \left(
        \frac{
        Q_{\beta_k}(\{g_j(r)\},\{w_j\})
        }{2}
    \right).
\]
Each $h_k$ is clamped to $[L,U]$. Duplicate rounded gate counts are removed.
If this leaves fewer than $K$ unique probes, the remaining probes are filled
from linearly spaced even gate counts across $[L,U]$.

\section{Batch Submission and Local Brackets}

For each ratio, the selected gate counts are converted to RMB configurations
and submitted as a stitched batch. The initial ratio uses
\texttt{initial\_batch\_shots} shots per point. Later ratios use
\texttt{trace\_batch\_shots} shots per point. The batch submitter splits
requests when necessary so that each stitched submission respects
\texttt{max\_cost\_per\_run}. The global budget and the optional per-ratio
budget are charged by the same HQC accounting used by the other RMB
experiments.

For reporting a local crossing, the method scans the probed gate counts in
ascending order. Let the Bayesian posterior probability that a measured
configuration is above the threshold be
\[
    a(n,r)=\Pr(p(n,r)>0.5\mid \text{local Beta posterior}).
\]
The helper declares a point confidently above when
\[
    a(n,r)\ge \texttt{decision\_confidence},
\]
and confidently below when
\[
    1-a(n,r)\ge \texttt{decision\_confidence}.
\]
The first adjacent transition from a confidently above point to a confidently
below point gives a bracket
\[
    [n_{\mathrm{lo}}, n_{\mathrm{hi}}].
\]
The recorded crossing is the even-rounded midpoint:
\[
    n_{\mathrm{cross}}
    =
    2\operatorname{round}
    \left(
        \frac{(n_{\mathrm{lo}}+n_{\mathrm{hi}})/2}{2}
    \right).
\]
This crossing is stored and plotted, but it is not the object used to update
the posterior. The posterior update uses the raw Bernoulli counts at all
measured configurations.

\section{Per-Ratio Budget}

The first batch may spend from the full remaining global budget. After at
least one measured configuration exists, the method optionally caps each ratio
by
\[
    B_r =
    \min(\texttt{max\_hqc\_per\_ratio}, B),
\]
where $B$ is the remaining global budget. If
\texttt{max\_hqc\_per\_ratio} is \texttt{None}, the ratio uses the full
remaining global budget. A batch is only submitted if its stitched HQC cost is
affordable under both the global and ratio budgets.

\section{Complete Algorithm}

\begin{enumerate}[leftmargin=*]
    \item Initialize RNG, RMB, backend, data store, and global budget with
    \texttt{start\_run(settings)}.
    \item Set the ratio sequence from $r_{\max}$ down to $r_{\min}$, using
    either fixed or adaptive ratio steps.
    \item For the first ratio:
    \begin{enumerate}
        \item Build the initial analytic prior window using the configured
        initial error guesses and uncertainty.
        \item Probe \texttt{initial\_batch\_points} gate counts with
        \texttt{initial\_batch\_shots} shots each.
        \item Record a local crossing if the batch contains a confident
        above-to-below transition.
    \end{enumerate}
    \item For each later ratio $r$:
    \begin{enumerate}
        \item If fewer than \texttt{posterior\_min\_configs} measured
        configurations exist, fall back to the initial analytic window.
        \item Otherwise sample prior particles
        $\theta_j=(q_j,s_j,c_j,\kappa_j)$.
        \item Compute likelihood weights from all measured Bernoulli counts.
        \item Evaluate each particle boundary
        $g_j(r)=1/(q_j+s_jr+c_jr(1-r))$.
        \item Use posterior weighted quantiles to build the next gate-count
        window.
        \item Choose probe gate counts from posterior weighted quantiles.
        \item Submit the batch with \texttt{trace\_batch\_shots} shots per
        point.
        \item Record a local crossing if the batch contains a confident
        above-to-below transition.
    \end{enumerate}
    \item Continue until the ratio sequence ends or the global budget is
    exhausted.
    \item Save the RMB data and the recorded crossings, print the experiment
    summary, and optionally plot the posterior fit and uncertainty bands.
\end{enumerate}

\section{Reproduction Notes}

To reproduce the method exactly, the following implementation details matter.

\begin{itemize}[leftmargin=*]
    \item Gate counts are rounded to even values before being converted to
    RMB configurations.
    \item The same measured data are reused at every posterior update; the
    posterior is not updated only from the most recent batch.
    \item Prior particles are resampled for each update from the same
    configured prior, then reweighted by the full likelihood. This is an
    importance-sampling approximation to the posterior.
    \item The posterior is not resampled sequentially between ratios. All
    adaptation enters through the likelihood of the accumulated data.
    \item Accepted crossings are diagnostic outputs. They do not define the
    posterior fit.
    \item The public method named \texttt{boundary\_bootstrap} returns
    posterior particle draws for compatibility with shared plotting code.
    \item The default local backend is a SympleQ backend with one-qubit Pauli
    error $2.5\times 10^{-5}$ and two-qubit Pauli error $7.9\times 10^{-4}$.
\end{itemize}

\end{document}
